\documentclass[11pt]{article}
\usepackage[margin=1in]{geometry}
\usepackage{amsmath,amssymb,booktabs,array}
\usepackage[round,authoryear]{natbib}
\usepackage[colorlinks=true,linkcolor=blue,citecolor=blue,urlcolor=blue]{hyperref}
\setlength{\parskip}{0.45em}
\setlength{\parindent}{0pt}
\newcommand{\E}{\mathbb{E}}
\newcommand{\A}{\mathcal{A}}
\newcommand{\Pg}{\mathfrak{P}}
\newcommand{\code}[1]{\texttt{\small #1}}

\title{Deposit funding and cross-border sovereign portfolios\\[2pt]
\large Diagnosis, tractable alternatives, and a stress test for the global model}
\author{Working note, branch \code{analysis} (from \code{main} at \code{55e174d})}
\date{12 October 2026}

\begin{document}
\sloppy
\maketitle

\section*{Summary}

\textbf{Issue 1 (deposit funding rate).} In every deposit closure the code offers, the deposit
rate is an \emph{equilibrium price that falls} when sovereign risk rises, and that fall cancels
51--61\% of the credit-spread increase before it reaches firms' working-capital cost. The
model's own output decomposition attributes $-0.179\%$ of impact output to the credit spread
and $+0.104\%$ to the deposit rate, for a total of $-0.073\%$. The credit-spread channel alone
is close to Bocola's open-economy benchmark of $-0.157\%$; the funding-rate offset is what
separates the model from it. The two closures also cost tractability in different ways. The
national closure needs three solver devices and cannot host the TPI. The union closure needs an
extra state whose stationarity rests on a borrowed premium. And the choice between the two
flips the sign of the 25-quarter cumulative output response ($-3.9$ against $+0.5$). In the
stress test, the union closure's uncalibrated premium $\kappa_{\rm nfa}$ moves impact output four
times as much per log unit as the portfolio cost $\psi$.

\textbf{Issue 2 (cross-border holdings).} The adjustment cost $\psi$ does four jobs at once: it
selects a portfolio that single-$\lambda$ banks leave indeterminate; it sets the return
semi-elasticity of cross-border holdings, which is exactly $1/\psi$; under national funding
it is the inverse elasticity of the country's \emph{entire} external adjustment; and it keeps
the default branch solvable, since below $\psi\approx1$ the union economy has no solution on
the current ladder. Over the solvable range ($\psi\in[1,32]$, quick grid) impact output moves
by a third of its value, while the deposit-rate channel is about five times larger. Two calibration
faults compound this. The anchor is a \emph{level}, so a growing debt stock is absorbed at home
by construction. And the symmetric gross position gives the D bank a foreign share of its
sovereign book of 20\%, against 0.75\% in the EBA 2010 data.

\textbf{Recommendations.} (i) Fund the marginal unit of bank liabilities through a Eurosystem
refinancing facility at a policy rate, with sovereign collateral and a penalty schedule
(option D3). It nests Bocola's own open-economy closure, puts the baseline and the TPI on one
funding structure, and makes the deposit rate move only when collateral binds. (ii) Anchor
cross-border holdings on \emph{shares} with a return semi-elasticity calibrated to an observable
(option P1). Move to fixed shares (P3, two fewer states) once the funding side carries the
external margin, if the stress test confirms that the portfolio channel is then small.
(iii) Finish the stress test of Section~\ref{sec:stress}: the script is provided, the
quick-grid union sweep is done, and the national sweep and the full-grid confirmation remain.

\section{The two specifications, and what is wrong with them}
\label{sec:diag}

Notation follows the paper draft. Countries $X\in\{D,F\}$; $\varkappa=\mathcal M_F/\mathcal M_D=8$
is the size of $F$ relative to $D$; $p_t$ is the price of the $F$ good in $D$ goods.
$R_{X,t}\equiv1+r_{X,t}$ is the deposit rate set at $t$ and paid at $t+1$. $\Omega_{X,t,t+1}$ is
the banker's stochastic discount factor (Bocola's kernel). $\lambda_X$ is the divertible fraction
of assets $\A_{X,t}$, $n_{X,t}$ is net worth, and $\mu_{X,t}$ is the multiplier on the incentive
constraint $\alpha_{X,t}n_{X,t}\ge\lambda_X\A_{X,t}$, with Bocola's closed form
\begin{equation}
\mu_{X,t}=\max\Bigl\{1-\frac{\E_t[\Omega_{X,t,t+1}]R_{X,t}\,n_{X,t}}{\lambda_X\A_{X,t}},\,0\Bigr\},
\qquad
r^{wc}_{X,t}=r_{X,t}+\frac{\lambda_X\mu_{X,t}}{\E_t[\Omega_{X,t,t+1}]}.
\label{eq:mu}
\end{equation}
The loan rate $r^{wc}$ (the rate on working-capital loans) is the only channel from the financial
block to output, through labour demand $w_{X,t}/(1+\zeta r^{wc}_{X,t})$, where $\zeta$ is the
working-capital share.

\subsection{Deposit funding}

\textbf{National market} (the default on the unmerged branch \code{domestic-funding-fix}):
$\mathrm{dep}_{X,t}=P^c_{X,t}HA_{X,t}$ for each $X$, where $HA$ is household deposits in basket
units and $P^c$ the consumption price index; $R_{X,t}$ is whatever clears it.

\textbf{Union market} (\code{main}, and the TPI's economy on every branch): one union-wide
clearing condition plus real interest parity with the debt-elastic premium of
\citet{schmittgrohe2003closing},
\begin{gather}
\mathrm{dep}_{D,t}+\varkappa p_t\,\mathrm{dep}_{F,t}=P^c_{D,t}HA_{D,t}+\varkappa p_tP^c_{F,t}HA_{F,t},
\label{eq:union-clear}\\
R_{D,t}=R_{F,t}\frac{\E_t[p_{t+1}]}{p_t}-\kappa_{\rm nfa}\frac{V^{\rm flow}_{D,t}}{\bar Y_D},
\qquad V_{t+1}=R_{D,t}V^{\rm flow}_{D,t},
\label{eq:uip}
\end{gather}
where $V^{\rm flow}_{D,t}$ is $D$'s net cross-border deposit position and $\kappa_{\rm nfa}=0.01$
(\code{point\_map.py:451}, \code{calibration/global\_projection.py:46}).

\paragraph{Fault 1: the funding rate is an automatic stabiliser with the wrong sign.}
In both closures the deposit rate is the price of a riskless claim whose demand by banks falls
when sovereign risk rises (banks shrink and the bond losses hit net worth). The rate therefore
falls, and because $r^{wc}=r+\text{spread}$, the fall offsets the spread one for one:

\begin{center}
\small
\begin{tabular}{lccc}
\toprule
Headline shock ($p^d=2\%$/qtr), impact & union, \code{main} & union, current & national, current\\
\midrule
credit spread $\lambda\mu/\E[\Omega]$ (bp/yr) & $+79.4$ & $+76.6$ & $+65.0$\\
deposit rate $r_D$ (bp/yr) & $-47.9$ & $-46.6$ & $-33.1$\\
net loan rate $r^{wc}_D$ (bp/yr) & $+34.5$ & $+32.9$ & $+34.1$\\
share of the spread cancelled & 60\% & 61\% & 51\%\\
\addlinespace
output, credit-spread channel (\%) & $-0.179$ & $-0.173$ & $-0.146$\\
output, deposit-rate channel (\%) & $+0.104$ & $+0.101$ & $+0.072$\\
output, total impact (\%) & $-0.073$ & $-0.061$ & $-0.068$\\
\bottomrule
\end{tabular}
\end{center}
{\footnotesize Full grid (\code{S\_REFINE}=5). ``Current'' is \code{domestic-funding-fix} with
the household-budget and cross-border $p'/p$ fixes. Source: \code{results/GLOBAL},
\code{results/UNION\_current}, \code{results/NATIONAL} run logs, \code{output\_decomposition}.}

The credit-spread channel alone ($-0.15$ to $-0.18\%$) lands on Bocola's like-for-like target
for his open economy, $-0.157\%$, whose transmission (GHH preferences plus working capital) this
model shares. In that economy the funding rate is a world rate,
$R=1/\beta+0.01\,B^{\rm for}/\mathrm{gdp}$, which hardly moves on impact. The gap between our
impact response and his is therefore, to first order, the endogenous funding rate. (This is an
accounting decomposition, not a counterfactual: a closure that pins the rate also changes the
general equilibrium. The stress test can measure that directly.)

The fall is also counterfactual. In 2010--12 periphery bank funding costs \emph{rose}: deposit
rates in Greece rose relative to the euro area, private deposits ran off, and the Eurosystem
(MROs, LTROs, ELA) became the marginal funder. Bank and sovereign credit risk co-moved
\citep{acharya2014pyrrhic}, and sovereign stress reached bank lending through banks' sovereign exposures
\citep{altavilla2017bank}. The model cannot produce this, because deposits are riskless in both
regimes (the household-side ``$\pi$-blindness'' listed in \code{CLAUDE.md}) and nothing else
connects the deposit rate to the sovereign.

\paragraph{Fault 2: each closure buys solvability with a device.}
\emph{National:} the D market has no external margin except the banks' cross-border sovereign
books. Solving it at the 55\% haircut needs no-short-sale Fischer--Burmeister pairs on all three
bank books, a collapsed TPI band, and a least-squares fallback in the Newton. A purchase barely
moves the D bond price, so the TPI cap cannot be held below about 490 bp, and the backstop is
evaluated on a different economy (union) from the paper's baseline (national).
\emph{Union:} the external position $V$ is a twelfth state whose only anchor is $\kappa_{\rm nfa}$,
a stationarity device borrowed from the small-open-economy literature and never calibrated here.
Parity also puts an expectation $\E_t[p_{t+1}]$ inside a residual.

\paragraph{Fault 3: the closure, not the friction, decides the medium run.}
On the current code, cumulative output over 25 quarters is $-3.89$ (national) against $+0.51$
(union); cumulative investment is $-6.00$ against $+4.07$. The impact responses are almost the
same. A reader cannot tell from the paper which of these is the model's answer.

\subsection{Cross-border holdings}

The F bank's first-order condition for the D bond, with its non-pecuniary cost, is
(\code{point\_map.py:356--357, 465}; the $p_t/p_{t+1}$ factor is in the draft and on
\code{domestic-funding-fix}, \emph{not on \code{main}})
\begin{equation}
\E_t\Bigl[\Omega_{F,t,t+1}\Xi^D_{t+1}\frac{p_t}{p_{t+1}}\Bigr]
=q^D_t\bigl(\E_t[\Omega_{F,t,t+1}]R_{F,t}+\lambda_F\mu_{F,t}\bigr)
\Bigl[1+\psi_F\frac{b^D_{F,t+1}-\bar b^D_F}{\bar B_D}\Bigr],
\label{eq:psi}
\end{equation}
with $\psi_F=\psi_D=2$ and $\bar b^D_F=\bar b^F_D=0.196$ (per $D$ capita). Here $\Xi^D$ is the
bond's payoff (coupon plus continuation value, times recovery in default), $q^D$ its price,
$b^D_F$ the F bank's holding, and $\bar B_D$ the SS debt stock.

\paragraph{What $\psi$ is.} Define the F bank's expected excess return on the D bond, net of its
marginal funding cost including the shadow cost of its constraint,
\begin{equation}
e^D_{F,t}\equiv\frac{\E_t[\Omega_{F,t,t+1}\Xi^D_{t+1}p_t/p_{t+1}]}
{q^D_t(\E_t[\Omega_{F,t,t+1}]R_{F,t}+\lambda_F\mu_{F,t})}-1 .
\end{equation}
Then \eqref{eq:psi} is \emph{exactly}
\begin{equation}
\frac{b^D_{F,t+1}-\bar b^D_F}{\bar B_D}=\frac{e^D_{F,t}}{\psi_F}.
\label{eq:semi}
\end{equation}
So $1/\psi$ is the semi-elasticity of cross-border holdings, as a fraction of the stock, with
respect to the quarterly excess return. At $\psi=2$, an excess return of 100 bp/yr moves holdings
by 0.125\% of the stock. Because the D bank prices the bond with no friction, $e^D_F$ is the gap
between the two banks' risk-adjusted valuations. That gap is driven by how differently
$\Omega_D$ and $\Omega_F$ load on the default state, and by $\mu_D$ against $\mu_F$.

\paragraph{Four jobs for one parameter.} (a) \emph{Selection.} A single-$\lambda$ banker's
value is linear in its portfolio, so at a symmetric SS the split of each bond between the two
banks is indeterminate, and off the SS it jumps to a corner. $\psi>0$ is what makes it
determinate. (b) \emph{Elasticity.} By \eqref{eq:semi}, $\psi$ alone scales the retrenchment
or flight-home response. (c) \emph{External margin.} Under national funding, $D$'s external
position is the net cross-border sovereign book
$\mathrm{nfa}_D=p\,q^F b^F_D-q^D b^D_F$. Every unit of the current account therefore passes
through \eqref{eq:psi} and its twin, and $\psi$ becomes the inverse elasticity of $D$'s access to
foreign saving. That is the mechanism behind the national ``$S=I$ trap'' (Fault~3), and it is
why one should expect $\psi$ to matter far more under national than under union funding.
(d) \emph{Solvability.} In the default branch, at the 55\% haircut, F banks want to sell their
entire D book. A finite $\psi$ is what keeps that branch interior: the sweep finds no solution
at $\psi\le0.5$ (Section~\ref{sec:stress}), the same corner that forced the no-short-sale pairs
under national funding.

\paragraph{Calibration faults.} (e) \emph{The anchor is a level.} When the risk shock raises the
debt stock, \eqref{eq:psi} makes F banks resist holding more than $\bar b^D_F$, so new issuance is
absorbed at home and the F share falls mechanically. That is home bias produced by the units, not
by a mechanism. The quick-grid sweep below shows it: a larger $\psi$ gives a \emph{larger} fall
in the F share. (f) \emph{The SS position contradicts the data on the D side.} The symmetric gross
positions required for a zero net position at $p=1$ give:

\begin{center}
\small
\begin{tabular}{lcc}
\toprule
& model SS & EBA 2011 stress test (Dec 2010)\\
\midrule
F-bank share of bank-held D debt & 20.0\% & 12.7\%\\
foreign share of the D bank's sovereign book & 20.0\% & 0.75\%\\
foreign share of the F bank's sovereign book & 2.5\% & 2.45\%\\
\bottomrule
\end{tabular}
\end{center}
{\footnotesize EBA: Greek banks hold EUR 54.4bn of Greek and 0.41bn of German sovereigns; German
banks hold 7.93bn of Greek and 315.3bn of German (\code{data/eba\_moments.json}, gross direct long).}

So the D bank carries a buffer of foreign safe bonds about 27 times the one Greek banks
actually had. Under national funding that buffer is $D$'s only way to finance a current-account
gap.

\section{Literature}
\label{sec:lit}

\paragraph{Funding rates in intermediary models.} In \citet{gertler2011model} and
\citet{gertler2010financial} deposits are riskless and the rate comes from the household Euler
equation. \citet{gertler2010financial} add an interbank market in which a fraction $\omega$ of
interbank borrowing escapes the divertibility problem, which spans segmented and frictionless
funding. \citet{bocola2016} keeps deposits riskless; his open-economy section (§V.C) closes the
model with a world rate plus a debt-elastic premium in the manner of
\citet{schmittgrohe2003closing}. In emerging-market models with banks,
\citet{aoki2016monetary} let banks borrow both domestically and abroad, with foreign funding more
easily diverted; the convex divertibility $\theta(1+\tfrac{\gamma}{2}x^2)$, where $x$ is the
foreign share, prices the funding mix. \citet{akinci2018exchange} and
\citet{dedola2013global} build two-country intermediary models with cross-border funding and
holdings; \citet{kollmann2011global} model a global bank. Country-spread closures, in which the
borrowing rate is the world rate times a country premium, are standard since
\citet{neumeyer2005business} and \citet{uribe2006country}. In \citet{neumeyer2005business} the
spread works through working capital, which is this model's own output channel. On the
sovereign-to-private pass-through, \citet{corsetti2013sovereign} posit that private funding
costs load a fraction of the sovereign spread. Deposit-supply frictions and convenience yields
\citep{drechsler2017deposits,begenau2020capital} make deposit rates sticky relative to the
safe rate. On the euro-area plumbing, \citet{drechsler2016who} show that weakly capitalised
banks borrowed most from the lender of last resort against risky collateral, and
\citet{crosignani2020unintended} show that LTRO funding was used to buy domestic sovereigns.
\citet{gertler2016wholesale} split bank funding into retail and wholesale and add runs.

\paragraph{Cross-border portfolios.} \citet{devereux2011country} and \citet{tille2010international}
determine portfolios from second- and third-order terms around a deterministic SS. That is not
needed in a global solution, but they show why first-order models need an extra device.
Portfolio-balance and limited-risk-bearing closures, from \citet{blanchard2005international} to
\citet{gabaix2015international} and \citet{itskhoki2021exchange}, make holdings a finite-elastic
function of expected excess returns, with the elasticity tied to intermediaries' risk capacity.
\citet{coeurdacier2013home} survey home bias. Demand-system estimation
\citep{koijen2019demand,koijen2021inspecting} measures exactly the object in \eqref{eq:semi},
the elasticity of a sector's holdings to yields, for euro-area sovereigns by holder sector. For
the euro crisis, \citet{broner2014sovereign} and \citet{brutti2016repatriation} document the
repatriation of periphery debt to domestic holders. Explanations include creditor
discrimination \citep{broner2014sovereign}, risk-shifting or carry trades by weak banks
\citep{acharya2015greatest}, moral suasion \citep{ongena2019invisible}, financial repression
\citep{becker2018financial}, and the bank--sovereign doom loop
\citep{farhi2018deadly,gennaioli2014sovereign}. \citet{altavilla2017bank} and
\citet{battistini2014systemic} measure the exposures and their transmission. Value-at-risk
constraints \citep{adrian2014procyclical} give a structural reason for risk-sensitive exposure
limits.

\section{Alternatives for the deposit funding rate}
\label{sec:dep}

Each option is stated with what it does to the period map (\code{point\_map.py}): its
residuals, unknowns and states.

\subsection*{D1. Two-tier funding with costly cross-border wholesale funding}

The D bank funds itself with net worth, national retail deposits $\mathrm{dep}_D$ at $R_D$, and
union wholesale funding $w_D$ (D goods) at $R^w$. Let $x_t\equiv w_{D,t}/\A_{D,t}$. In the spirit
of \citet{aoki2016monetary}, wholesale funding raises divertibility:
\begin{equation}
\alpha_t n_t\ \ge\ \lambda\,\Theta(x_t)\A_t,\qquad \Theta(x)=1+\tfrac{\gamma}{2}x^2 .
\end{equation}
Because $\Theta(x)\A$ is homogeneous of degree one in $(\A,w)$, Bocola's closed form survives
with $\lambda\to\lambda\Theta(x)$:
\begin{gather}
\mu_t=\max\Bigl\{1-\frac{\E_t[\Omega]R_{D,t}\,n_t}{\lambda\Theta(x_t)\A_t},0\Bigr\},\qquad
\alpha_t=\frac{\E_t[\Omega]R_{D,t}}{1-\mu_t},\\
\E_t[\Omega(R^K_{t+1}-R_{D,t})]=\lambda\mu_t\bigl(1-\tfrac{\gamma}{2}x_t^2\bigr)
\quad\text{(and likewise for bonds and loans)},\\
R_{D,t}-R^w_t=\frac{\lambda\gamma\,\mu_t\,x_t}{\E_t[\Omega]}+\kappa_w(x_t-\bar x),
\label{eq:d1-gap}
\end{gather}
with retail clearing $\mathrm{dep}_{D,t}=P^c_{D,t}HA_{D,t}$, parity for the union lender
$R^w_t=R_{F,t}\E_t[p_{t+1}]/p_t$, and $F$'s clearing net of its households' wholesale lending.
The small $\kappa_w$ term is an SGU device that matters only where $\mu=0$.
\emph{Nesting:} $\gamma\to\infty$ gives $x\to0$ (national); $\gamma=0$ gives $R_D=R^w$ (union).
\emph{Sign:} when $\mu_D$ rises, \eqref{eq:d1-gap} pushes the domestic deposit rate \emph{up}
relative to the union rate, the deposit competition of 2010--12.
\emph{Calibration in closed form:} at the SS, dividing \eqref{eq:d1-gap} by the credit spread
gives $\gamma=(R_D-R^w)/[\,\bar x\,(r^{wc}-r_D)\,]$ up to the $x^2$ term. That uses three
observables: the retail deposit spread over Euribor (ECB MIR statistics), the lending spread
(already targeted), and the non-resident wholesale share of Greek bank funding (Bank of Greece
MFI balance sheets).
\emph{Cost:} one unknown and one residual more (15); $V$ is reused as the wholesale position.

\subsection*{D2. Defaultable deposits: sovereign risk in the deposit rate}

Deposits pay a fraction $h^{\rm dep}\in(0,1]$ of their promised return in the default regime
(bail-in, as in Cyprus in 2013, or redenomination; \citealp{krishnamurthy2018ecb} measure the
latter). The household Euler weights the default branch:
\begin{equation}
x_{D,t}^{-\sigma}=\beta\,\E_t\Bigl[\bigl(1-(1-h^{\rm dep})d_{t+1}\bigr)R_{D,t}
\frac{P^c_{D,t}}{P^c_{D,t+1}}\,x_{D,t+1}^{-\sigma}\Bigr],
\label{eq:d2}
\end{equation}
where $x$ is the GHH composite and $d_{t+1}$ the default indicator. The bank's deposit leg is
written down by the same factor in that branch,
$\Pg^{d'=1}_{D,t+1}=h^{\rm dep}R_{D,t}\mathrm{dep}_{D,t}-R^{wc}_{D,t}L_{D,t}$,
so the depositor's loss is the bank's gain and nothing leaks from the budget. To first order,
\begin{equation}
r_{D,t}-r^{h=1}_{D,t}\ \approx\ \pi_t(1-h^{\rm dep})\,
\frac{\E_t[x^{-\sigma}_{t+1}\mid d'=1]}{\E_t[x^{-\sigma}_{t+1}]},
\end{equation}
so the deposit spread is the default probability times the deposit loss-given-default, weighted
by marginal utility. Matching the reduced form of \citet{corsetti2013sovereign}, deposit spread
$=\xi\times$ sovereign spread, gives $h^{\rm dep}\approx1-\xi(1-0.45)$.
\emph{Cost:} no new state and no new unknown. The default branch already exists in the
quadrature; only the next-period state for $d'=1$ becomes regime-specific inside
\code{\_continuation}, which builds the state per regime anyway. \emph{Caveat:} the bail-in
transfer cushions $n_D$ in default, so it lowers $\Omega^d$ and the bank-side risk premium;
redenomination without a bank-side gain is the other polar case. Works with any deposit market.

\subsection*{D3. Eurosystem refinancing at a policy rate, against sovereign collateral}

The D bank's marginal funding is a Eurosystem facility $\ell_t$ (D goods). The facility is
priced on a penalty schedule (ELA lends above the main refinancing rate) and limited by
haircut-adjusted collateral:
\begin{gather}
\A_{D,t}=n_{D,t}+\mathrm{dep}_{D,t}+\ell_t,\qquad
R^\ell_t=\bar R+\kappa_\ell\frac{\ell_t-\bar\ell}{\bar Y_D},\qquad
\ell_t\le(1-h^c)\,q^D_tb^D_{D,t+1},\\
R_{D,t}=R^\ell_t+\frac{\eta_t}{\E_t[\Omega_D]},\qquad
\eta_t\ge0,\quad (1-h^c)q^D_tb^D_{D,t+1}-\ell_t\ge0,\quad \eta_t\,\bigl[(1-h^c)q^D_tb^D_{D,t+1}-\ell_t\bigr]=0,
\label{eq:d3}\\
q^D_t\bigl(\E_t[\Omega_D]R_{D,t}+\lambda\mu_{D,t}-(1-h^c)\eta_t\bigr)=\E_t[\Omega_D\Xi^D_{t+1}].
\label{eq:d3-bond}
\end{gather}
Here $\eta$ is the collateral multiplier (in the bank's value units), $h^c$ the collateral
haircut, and $\kappa_\ell$ the slope of the penalty schedule. The Eurosystem funds $\ell$ with
reserves held by $F$ banks at $R_F$ and remits the facility's P\&L, including the
terms-of-trade valuation, by the capital key, through the same accounts the TPI already uses
($\Pi^{cb}$, \code{tpi\_key\_D}). That is what makes a rate fixed in D goods consistent in a
two-country model: the valuation profit is assigned.

\emph{What it does.} While collateral is slack, $R_D=R^\ell$ is set by policy: the deposit rate
stops absorbing the shock, households are no longer forced to hold whatever deposits the bank
wants, and $D$ can draw on union saving through the facility (TARGET2 in all but name). When the
sovereign is marked down, collateral binds and $\eta>0$, which raises $R_D$: the pass-through has
the right sign and runs through $q^D$. By \eqref{eq:d3-bond}, binding collateral also raises the
D bank's valuation of its own sovereign, so repatriation becomes an outcome rather than a
by-product of units \citep{crosignani2020unintended,drechsler2016who}.
\emph{Nesting:} with collateral slack and $\bar\ell=0$, $R^\ell=\bar R+\kappa_\ell\,\ell/\bar Y_D$
is Bocola's open-economy closure ($\kappa_\ell=0.01$, with $\ell$ in the role of his foreign
debt); $\kappa_\ell\to\infty$ with $\bar\ell=0$ is the national market.
\emph{Calibration:} $h^c$ comes from the published Eurosystem haircut schedule for the rating
class; $\bar\ell$ from Greek banks' Eurosystem borrowing before the crisis; $\kappa_\ell$ from the
ELA--MRO spread.
\emph{Cost:} one state (the facility obligation, replacing the inert $V$ under national funding),
one unknown and one Fischer--Burmeister pair, a device the code already uses. The national
solver devices, the collapsed TPI band and the separate union baseline should no longer be
needed, because the external margin now exists in every economy.

\begin{center}
\small
\begin{tabular}{p{2.6cm}p{2.5cm}p{2.7cm}p{3.5cm}p{3.0cm}}
\toprule
& sign of $\Delta r_D$ on risk & nests & new parameters (data) & states / unknowns\\
\midrule
national (now) & falls & --- & --- & 12 / 14, 3 devices\\
union (now) & falls & --- & $\kappa_{\rm nfa}$ (none) & 12 / 14\\
D1 wholesale & rises if $\mu\uparrow$ & national, union & $\gamma,\bar x$ (MIR, BoG) & 12 / 15\\
D2 defaultable & rises with $\pi$ & any closure & $h^{\rm dep}$ (MIR spread) & unchanged\\
D3 Eurosystem & rises if collateral binds & national, Bocola §V.C & $h^c,\bar\ell,\kappa_\ell$ (ECB) & 12 / 15, 1 FB pair\\
\bottomrule
\end{tabular}
\end{center}

\textbf{Recommendation: D3.} It is the only option in which the marginal funding rate is a
policy rate rather than an equilibrium price, which removes Fault~1 by construction outside the
collateral-binding region. It removes Fault~2, because baseline and TPI share one funding
structure with the central bank in both. Its three parameters have published counterparts. D2
can be layered on top if the paper wants deposit flight. D1 is the cleanest dial for a
robustness section, since $\gamma$ moves the economy continuously from national to union.

\section{Alternatives for cross-border holdings}
\label{sec:port}

\subsection*{P1. Share anchor with an identified return semi-elasticity}

Replace the level anchor in \eqref{eq:psi} by the holder's share of the outstanding stock,
$s_{F,t}\equiv b^D_{F,t+1}/B^D_{t+1}$:
\begin{equation}
\E_t\Bigl[\Omega_F\Xi^D_{t+1}\frac{p_t}{p_{t+1}}\Bigr]
=q^D_t\bigl(\E_t[\Omega_F]R_{F,t}+\lambda\mu_{F,t}\bigr)\bigl[1+\tilde\psi_F(s_{F,t}-\bar s_F)\bigr]
\iff s_{F,t}-\bar s_F=\frac{e^D_{F,t}}{\tilde\psi_F}.
\end{equation}
New issuance is then shared at $\bar s_F$ unless returns say otherwise. $1/\tilde\psi$ is the
share's semi-elasticity, which can be read from demand-system estimates for foreign banks
\citep{koijen2021inspecting}, or from the 2010--11 change in German banks' Greek exposure (EBA
2011 stress test against the EBA 2011 capital exercise; BIS consolidated banking statistics)
divided by the model-implied $\Delta e$ at the observed spread. Set $\bar s_F=12.7\%$ (EBA). Note
that matching the D side (0.75\%) breaks the zero net position at $p=1$. Under D3 the facility
can carry $D$'s SS net debt, so the asymmetry is absorbed on the funding side without moving
$p_{\rm ss}$. \emph{Cost:} none; same residuals and states.

\subsection*{P2. Limited risk-bearing capacity (a risk budget)}

GK kernels price $D$'s default as nearly idiosyncratic for an $F$ banker, so the model has no
force that makes foreign banks \emph{shed} when risk rises. A risk budget in the spirit of
\citet{adrian2014procyclical} and \citet{gabaix2015international} adds one. Let
$\chi_{F,t}\equiv q^D_tb^D_{F,t+1}/(\varkappa p_tn_{F,t})$ be the exposure as a share of $F$ net
worth, and let $\mathcal V_t\equiv\mathrm{Var}_t\bigl(\Xi^D_{t+1}p_t/(q^D_tp_{t+1})\bigr)$ be its
conditional return variance, computed on the existing $s'\times d'$ quadrature at no extra cost.
Then
\begin{equation}
\chi_{F,t}\,(\bar\sigma^2+\mathcal V_t)=\chi_{\rm ss}\,\bar\sigma^2+\frac{e^D_{F,t}}{\gamma_F}.
\label{eq:p2}
\end{equation}
At the deterministic SS ($\mathcal V=0$, $e=0$) $\chi=\chi_{\rm ss}$, so it is determinate. When
$\pi$ rises, $\mathcal V_t$ rises roughly in proportion to $\pi_t(1-0.45)^2$ and the exposure
falls by the factor $\bar\sigma^2/(\bar\sigma^2+\mathcal V_t)$; when $F$ net worth falls, the
level falls too, which carries contagion. $\bar\sigma^2$ is the variance at which the budget
halves the exposure; $\gamma_F$ is the return sensitivity. Both are identified by the
2010--11 retrenchment together with its return. \emph{Cost:} same residual slots, a few
quadrature sums.

\subsection*{P3. Fixed shares (segmentation, the $\psi\to\infty$ limit)}

$b^D_{F,t+1}=\bar s_FB^D_{t+1}$ and $b^F_{D,t+1}=\bar s_DB^F$ ($B^F$ is fixed in the code).
Residuals 9 and 13 and their two unknowns go, and the states $b_{DF}$ and $b_{FD}$ become
functions of $b_{DD}$ and $M^{cb}$, so the grid goes from 12 states to 10. On the measured
scaling (the solve is $\propto n^2$, with 115 points falling to 95), the full grid is about 30\%
cheaper. F banks still take marked-to-market losses on $\bar s_F$ of the D stock, so contagion
survives; what is lost is cross-bank rebalancing, and with it the split of the TPI's purchases
by seller. \emph{Only admissible with an external funding margin} (union, D1 or D3). Under
national funding it closes the economy to foreign saving altogether.

\textbf{Recommendation: P1 now; P3 once D3 is in, if Section~\ref{sec:stress} finds the
portfolio channel small there; P2 if retrenchment dynamics become a target.} The general point
is that the model needs exactly one external adjustment margin with a calibratable elasticity.
Today, depending on the deposit closure, that margin is either a frictionless deposit position
(union, $\kappa_{\rm nfa}$) or the $\psi$-governed bond books (national). D3 puts it on the
funding side, where 2010--12 actually put it, and so leaves $\psi$ with one job, not three.

\section{Stress test: how much does $\psi$ drive the dynamics?}
\label{sec:stress}

\paragraph{Design.} The question is quantitative and global, so each point is a full re-solve,
never a perturbation of fixed rules.
\begin{enumerate}
\item \emph{Sweep.} $\psi\in\{0.25,0.5,1,2,4,8,32\}$ on a log grid, with both $\psi_F$ and
$\psi_D$ moved together (then one at a time). Solve the calibrated point cold, then walk
outward, warm-starting each solve from its neighbour: a continuation in $\psi$ that keeps every
Newton inside its basin. Quick grid ($\mu=1$, 25 points) for the sweep; the full grid
(\code{S\_REFINE}=5) at $\{0.5,2,8\}$ to confirm that the elasticities are grid-independent.
\item \emph{Read.} Bocola's differenced IRF from each solution's own stochastic rest point,
with the headline shock $p^d=2\%$/qtr for 25 quarters. Re-locate the rest point for every
$\psi$: it moves with $\psi$, and differencing against a stale one charges the drift to the
shock.
\item \emph{Metrics.} Impact, trough and cumulative output; cumulative investment; impact
$n_D$, $n_F$ and $q^D$; the credit spread and deposit rate on impact; the F-bank share of $D$
debt at $q4$; $D$'s net external position at $q8$; and the model's retrenchment elasticity
(pp of $D$ debt shed by F banks per 100 bp of sovereign spread), which is the moment to set
against the data.
\item \emph{Sensitivity.} The log-elasticity at the calibrated point,
$\varepsilon_m=\partial\ln|m|/\partial\ln\psi$, from the two neighbours, and the sweep range
$(\max m-\min m)/|m(\psi_{\rm ref})|$. Use elasticities rather than raw changes so that dials
measured in different units can be ranked against each other (step 7).
\item \emph{Isolation.} Two decompositions. (a) \emph{Portfolio channel}:
$m(\psi)-m(\psi\to\infty)$, the IRF relative to fixed holdings (P3 is that limit). (b)
\emph{Within-period versus expectations}: re-clear the period map exactly at impact under each
$\psi$ while holding the reference continuation rules fixed. The ratio of that change to the
fully re-solved change is the share of $\psi$'s effect that runs through the current FOC; the
remainder runs through expectations of future rebalancing.
\item \emph{Accuracy guards.} Each point records the solver verdict and worst residual
(Bocola's $\sum F^2\le m\cdot10^{-18}$), the largest box clip along the IRF (the bond-state
band is $\pm12\%$, so a small $\psi$ can push holdings out of the box), and the impact read two
ways, fitted and exactly cleared. A $\psi$ effect smaller than that gap is not identified at
the current resolution. On the full grid, add the Euler-error report on the ergodic set.
\item \emph{Ranking against the other dials.} Repeat the sweep for $\kappa_{\rm nfa}$, $f$, the
credit-spread target, $\zeta_{wc}$, the Jermann $\xi$, $\sigma_s$ and the recovery rate, and rank
by $|\varepsilon_{\rm cumY}|$. That gives a tornado table, from one-at-a-time sweeps. For
interactions, a Morris elementary-effects screen \citep{morris1991factorial} with ten
trajectories over these eight dials is about 90 quick solves, roughly an overnight run at
about 1--5 minutes each.
\item \emph{Interaction with the funding closure.} The whole sweep in both deposit markets.
Under national funding, $\psi$ is predicted (Section~\ref{sec:diag}) to govern the medium run;
under union funding only the composition. The ratio of the two elasticities is the cleanest
test of that claim.
\end{enumerate}

\paragraph{Implementation.} \code{Claude files/diagnostics/portfolio\_stress/psi\_sweep.py} does
steps 1--6 for any calibration key, and step 7 as a ranking table (\code{--compare}) across saved
sweeps. Results go to \code{results/STRESS/<key>\_<market>\_s<refine>.json}, with a six-panel
figure. Typical calls:
\begin{verbatim}
S="Claude files/diagnostics/portfolio_stress/psi_sweep.py"
python3 "$S"                                   # psi, quick grid
python3 "$S" --values 0.5 8 --s-refine 5       # full-grid confirmation
python3 "$S" --param kappa_nfa --values 0.0025 0.04
python3 "$S" --set deposit_market=national     # domestic-funding-fix only
python3 "$S" --compare                         # rank all saved sweeps
\end{verbatim}
\code{--set deposit\_market=national} works only on \code{domestic-funding-fix}; on \code{main}
the key does not exist and the run is the union economy.

\subsection*{First results: quick grid, union deposits, \code{main}}

\begin{center}
\scriptsize
\setlength{\tabcolsep}{3.5pt}
\begin{tabular}{lcccccccccc}
\toprule
$\psi$ & $Y_0$ fit & $Y_0$ exact & $\sum Y$ & $\sum I$ & $n_{D,0}$ & spread$_0$ & $r_{D,0}$ & $\Delta s_F$ q4 & nfa q8 & $Y_0$ exact\\
 & \% & \% & 25q & 25q & \% & bp/yr & bp/yr & pp & \% $Y$ & $\psi{=}2$ rules\\
\midrule
1 & -0.047 & -0.060 & +1.64 & +8.24 & -5.09 & +94.8 & -47.8 & -0.342 & +0.701 & -0.069 \\
\textbf{2} & -0.064 & -0.076 & +1.63 & +8.49 & -5.20 & +100.1 & -50.6 & -0.445 & +0.756 & -0.076 \\
4 & -0.051 & -0.062 & +1.90 & +8.98 & -5.20 & +102.3 & -52.3 & -0.505 & +0.745 & -0.079 \\
8 & -0.047 & -0.056 & +2.03 & +9.19 & -5.22 & +103.7 & -53.2 & -0.537 & +0.744 & -0.081 \\
32 & -0.043 & -0.052 & +2.13 & +9.36 & -5.23 & +104.9 & -53.9 & -0.563 & +0.743 & -0.083 \\
\midrule
range / $|m(2)|$ & 33\% &  & 30\% & 13\% & 3\% & 10\% & 12\% & 50\% & 7\% &  \\
\bottomrule
\end{tabular}
\end{center}
{\footnotesize Headline shock $p^d=2\%$/qtr, differenced IRF from each solution's own rest point.
Coarse grid (25 points), a preview: at $\psi=2$ the full grid gives $-0.073\%$ against
$-0.064\%$ here. The last column re-clears the period map exactly at impact under each $\psi$ with the
$\psi=2$ continuation rules held fixed (the within-period effect alone). The range row is
$(\max m-\min m)/|m(\psi{=}2)|$ over the solved values. No solution at $\psi=0.5, 0.25$: the warm walk (1, 2 and 4 log-steps from $\psi=1$) and the cold ladder both fail.}

\begin{center}
\small
\begin{tabular}{lccccc}
\toprule
dial (values swept) & $\ln$ range & $Y_0$ range & $\sum Y$ range & $\Delta s_F$ range & nfa q8 range\\
\midrule
$\psi$ (1--32) & 3.47 & 33\% & 30\% & 50\% & 7\%\\
$\kappa_{\rm nfa}$ (0.0025--0.04) & 2.77 & 110\% & 56\% & 3\% & 9\%\\
\addlinespace
\multicolumn{6}{l}{per unit of $\ln$(dial): $Y_0$: $\psi$ 9.5\%, $\kappa_{\rm nfa}$ 40\%; \quad $\sum Y$: $\psi$ 8.7\%, $\kappa_{\rm nfa}$ 20\%}\\
\bottomrule
\end{tabular}
\end{center}
{\footnotesize Same grid, shock and reads as above; ranges as a share of the value at the
calibrated point. $\kappa_{\rm nfa}$: impact output $-0.085$, $-0.064$, $-0.014\%$ at
$0.0025$, $0.01$, $0.04$; cumulative output $+2.02$, $+1.63$, $+1.10$
(\code{results/STRESS/kappa\_nfa\_default\_s0.json}).}

\paragraph{Reading.}
\begin{enumerate}
\item \emph{Solvability bounds $\psi$ from below.} No solution exists at $\psi\le0.5$: the warm
walk from $\psi=1$ (in one, two and four log-steps) stalls at $\max|F|\approx3\cdot10^{-2}$, and
the cold ladder produces a non-finite iterate in the default-branch haircut homotopy. In that
branch the F bank wants to sell its whole D book; with the $10^{-4}$ floors in place of a
complementarity pair the Jacobian degenerates. This is the national branch's failure, reached
here through the portfolio cost. The interesting low-$\psi$ region cannot be explored without
porting the no-short-sale pairs to the union economy.
\item \emph{Under union funding, $\psi$ governs composition, not the current account.} The F
share moves by half its value across the sweep, while $D$'s net external position moves by 7\%,
because the deposit position $V$ absorbs the external adjustment whatever $\psi$ is (the
six-panel figure, \code{results/STRESS/psi\_bF\_D\_default\_s0.png}). This is the prediction of
Section~\ref{sec:diag}, job (c), read in reverse.
\item \emph{Output: a third over a 32-fold range, with a bump at the calibrated point.} A
stiffer friction gives a smaller impact fall and a larger later boom. Above $\psi=2$ both move
monotonically: the one-sided elasticity between 2 and 32 is $-0.15$ for $|Y_0|$ and $+0.10$
for cumulative output. The $\psi=1$ point breaks monotonicity on impact. That is \emph{not}
multiplicity: cold-solving $\psi=4$ reproduces the warm solution exactly ($-0.0514\%$), and
walking $4\to2$ warm reproduces the cold $\psi=2$ ($-0.0644\%$). On 25 points it may be an
interpolation effect at the KKT kink; the full-grid points $\{1,2,4\}$ decide.
\item \emph{The mechanism runs through expectations.} Holding the $\psi=2$ rules fixed and
changing $\psi$ only in the period map moves impact output the \emph{other} way ($-0.069\%$ at
$\psi=1$ to $-0.083\%$ at $\psi=32$): today a stiffer friction stops F banks buying at the dip,
so the D bank absorbs more. The fully re-solved effect has the opposite sign, so what $\psi$
does to the IRF is mostly what it does to expected future rebalancing, which a fixed-rule read
excludes.
\item \emph{The level anchor shows up as predicted.} The F share falls \emph{more} as $\psi$
rises ($-0.34$ to $-0.56$ pp at q4), the opposite of what a stiffer rebalancing friction should
do, because the friction holds the level while the debt stock grows. At $\psi=1$ the F share
rises on impact (F banks buy at the dip) before falling.
\item \emph{The funding closure dominates.} Per log unit, $\kappa_{\rm nfa}$ moves impact output
four times as much as $\psi$, and cumulative output twice as much. A stationarity device borrowed
from the small-open-economy literature, never calibrated here, is the most powerful dial of the
two. Its effect does not run through the loan rate (impact $r^{wc}$ is $+50$ bp/yr at
$\kappa=0.0025$ and $+48$ at $0.04$) but through the terms of trade and labour supply, as the
external margin tightens. The deposit-rate channel of the decomposition ($+0.104$ pp) is likewise
about five times the whole $\psi$ range of impact output (0.021 pp). For the \emph{union}
economy, then, the answer to Issue~2 is that $\psi$ drives the cross-border composition strongly
and the output path moderately, and Issue~1 is first-order. Under national funding the
prediction is the reverse; that sweep is the next run.
\end{enumerate}


\section{Next steps}

\begin{enumerate}
\item Merge or port the cross-border $p'/p$ fix to \code{main} before quoting any $\psi$ result;
the sweep above is on \code{main}'s FOC.
\item Run the sweep under national funding on \code{domestic-funding-fix}, and the full-grid
confirmation at $\{0.5,2,8\}$.
\item Pull the four data moments: the GR--euro area deposit spread (ECB MIR), the non-resident
wholesale share and Eurosystem borrowing of Greek banks (Bank of Greece, ECB), German banks'
Greek exposure in 2011 (EBA capital exercise, BIS CBS), and the ELA penalty spread.
\item Implement D3 behind a \code{deposit\_market="eurosystem"} switch (the N1--N5 nesting tests
extend directly: $\kappa_\ell\to\infty$ must reproduce national), then P1, re-run the sweep, and
decide on P3.
\end{enumerate}

\begin{thebibliography}{99}
\small
\bibitem[Acharya et~al.(2014)]{acharya2014pyrrhic} Acharya, V., I. Drechsler and P. Schnabl (2014). A pyrrhic victory? Bank bailouts and sovereign credit risk. \emph{Journal of Finance} 69(6).
\bibitem[Acharya and Steffen(2015)]{acharya2015greatest} Acharya, V. and S. Steffen (2015). The ``greatest'' carry trade ever? Understanding eurozone bank risks. \emph{Journal of Financial Economics} 115(2).
\bibitem[Adrian and Shin(2014)]{adrian2014procyclical} Adrian, T. and H. S. Shin (2014). Procyclical leverage and value-at-risk. \emph{Review of Financial Studies} 27(2).
\bibitem[Akinci and Queralto(2018)]{akinci2018exchange} Akinci, O. and A. Queralto (2018). Exchange rate dynamics and monetary spillovers with imperfect financial markets. Federal Reserve Bank of New York Staff Report.
\bibitem[Altavilla et~al.(2017)]{altavilla2017bank} Altavilla, C., M. Pagano and S. Simonelli (2017). Bank exposures and sovereign stress transmission. \emph{Review of Finance} 21(6).
\bibitem[Aoki et~al.(2016)]{aoki2016monetary} Aoki, K., G. Benigno and N. Kiyotaki (2016). Monetary and financial policies in emerging markets. Working paper.
\bibitem[Battistini et~al.(2014)]{battistini2014systemic} Battistini, N., M. Pagano and S. Simonelli (2014). Systemic risk, sovereign yields and bank exposures in the euro crisis. \emph{Economic Policy} 29(78).
\bibitem[Becker and Ivashina(2018)]{becker2018financial} Becker, B. and V. Ivashina (2018). Financial repression in the European sovereign debt crisis. \emph{Review of Finance} 22(1).
\bibitem[Begenau(2020)]{begenau2020capital} Begenau, J. (2020). Capital requirements, risk choice, and liquidity provision in a business-cycle model. \emph{Journal of Financial Economics} 136(2).
\bibitem[Blanchard et~al.(2005)]{blanchard2005international} Blanchard, O., F. Giavazzi and F. Sa (2005). International investors, the U.S. current account, and the dollar. \emph{Brookings Papers on Economic Activity} 2005(1).
\bibitem[Bocola(2016)]{bocola2016} Bocola, L. (2016). The pass-through of sovereign risk. \emph{Journal of Political Economy} 124(4).
\bibitem[Broner et~al.(2014)]{broner2014sovereign} Broner, F., A. Erce, A. Martin and J. Ventura (2014). Sovereign debt markets in turbulent times: Creditor discrimination and crowding-out effects. \emph{Journal of Monetary Economics} 61.
\bibitem[Brutti and Saur\'e(2016)]{brutti2016repatriation} Brutti, F. and P. Saur\'e (2016). Repatriation of debt in the euro crisis. \emph{Journal of the European Economic Association} 14(1).
\bibitem[Coeurdacier and Rey(2013)]{coeurdacier2013home} Coeurdacier, N. and H. Rey (2013). Home bias in open economy financial macroeconomics. \emph{Journal of Economic Literature} 51(1).
\bibitem[Corsetti et~al.(2013)]{corsetti2013sovereign} Corsetti, G., K. Kuester, A. Meier and G. M\"uller (2013). Sovereign risk, fiscal policy, and macroeconomic stability. \emph{Economic Journal} 123(566).
\bibitem[Crosignani et~al.(2020)]{crosignani2020unintended} Crosignani, M., M. Faria-e-Castro and L. Fonseca (2020). The (unintended?) consequences of the largest liquidity injection ever. \emph{Journal of Monetary Economics} 112.
\bibitem[Dedola et~al.(2013)]{dedola2013global} Dedola, L., P. Karadi and G. Lombardo (2013). Global implications of national unconventional policies. \emph{Journal of Monetary Economics} 60(1).
\bibitem[Devereux and Sutherland(2011)]{devereux2011country} Devereux, M. and A. Sutherland (2011). Country portfolios in open economy macro-models. \emph{Journal of the European Economic Association} 9(2).
\bibitem[Drechsler et~al.(2016)]{drechsler2016who} Drechsler, I., T. Drechsel, D. Marques-Ibanez and P. Schnabl (2016). Who borrows from the lender of last resort? \emph{Journal of Finance} 71(5).
\bibitem[Drechsler et~al.(2017)]{drechsler2017deposits} Drechsler, I., A. Savov and P. Schnabl (2017). The deposits channel of monetary policy. \emph{Quarterly Journal of Economics} 132(4).
\bibitem[Farhi and Tirole(2018)]{farhi2018deadly} Farhi, E. and J. Tirole (2018). Deadly embrace: Sovereign and financial balance sheets doom loops. \emph{Review of Economic Studies} 85(3).
\bibitem[Gabaix and Maggiori(2015)]{gabaix2015international} Gabaix, X. and M. Maggiori (2015). International liquidity and exchange rate dynamics. \emph{Quarterly Journal of Economics} 130(3).
\bibitem[Gennaioli et~al.(2014)]{gennaioli2014sovereign} Gennaioli, N., A. Martin and S. Rossi (2014). Sovereign default, domestic banks, and financial institutions. \emph{Journal of Finance} 69(2).
\bibitem[Gertler and Karadi(2011)]{gertler2011model} Gertler, M. and P. Karadi (2011). A model of unconventional monetary policy. \emph{Journal of Monetary Economics} 58(1).
\bibitem[Gertler and Kiyotaki(2010)]{gertler2010financial} Gertler, M. and N. Kiyotaki (2010). Financial intermediation and credit policy in business cycle analysis. \emph{Handbook of Monetary Economics} 3A.
\bibitem[Gertler et~al.(2016)]{gertler2016wholesale} Gertler, M., N. Kiyotaki and A. Prestipino (2016). Wholesale banking and bank runs in macroeconomic modeling of financial crises. \emph{Handbook of Macroeconomics} 2.
\bibitem[Itskhoki and Mukhin(2021)]{itskhoki2021exchange} Itskhoki, O. and D. Mukhin (2021). Exchange rate disconnect in general equilibrium. \emph{Journal of Political Economy} 129(8).
\bibitem[Koijen et~al.(2021)]{koijen2021inspecting} Koijen, R., F. Koulischer, B. Nguyen and M. Yogo (2021). Inspecting the mechanism of quantitative easing in the euro area. \emph{Journal of Financial Economics} 140(1).
\bibitem[Koijen and Yogo(2019)]{koijen2019demand} Koijen, R. and M. Yogo (2019). A demand system approach to asset pricing. \emph{Journal of Political Economy} 127(4).
\bibitem[Kollmann et~al.(2011)]{kollmann2011global} Kollmann, R., Z. Enders and G. M\"uller (2011). Global banking and international business cycles. \emph{European Economic Review} 55(3).
\bibitem[Krishnamurthy et~al.(2018)]{krishnamurthy2018ecb} Krishnamurthy, A., S. Nagel and A. Vissing-Jorgensen (2018). ECB policies involving government bond purchases: Impact and channels. \emph{Review of Finance} 22(1).
\bibitem[Morris(1991)]{morris1991factorial} Morris, M. (1991). Factorial sampling plans for preliminary computational experiments. \emph{Technometrics} 33(2).
\bibitem[Neumeyer and Perri(2005)]{neumeyer2005business} Neumeyer, P. A. and F. Perri (2005). Business cycles in emerging economies: The role of interest rates. \emph{Journal of Monetary Economics} 52(2).
\bibitem[Ongena et~al.(2019)]{ongena2019invisible} Ongena, S., A. Popov and N. Van Horen (2019). The invisible hand of the government: Moral suasion during the European sovereign debt crisis. \emph{American Economic Journal: Macroeconomics} 11(4).
\bibitem[Schmitt-Groh\'e and Uribe(2003)]{schmittgrohe2003closing} Schmitt-Groh\'e, S. and M. Uribe (2003). Closing small open economy models. \emph{Journal of International Economics} 61(1).
\bibitem[Tille and van Wincoop(2010)]{tille2010international} Tille, C. and E. van Wincoop (2010). International capital flows. \emph{Journal of International Economics} 80(2).
\bibitem[Uribe and Yue(2006)]{uribe2006country} Uribe, M. and V. Yue (2006). Country spreads and emerging countries: Who drives whom? \emph{Journal of International Economics} 69(1).
\end{thebibliography}

\end{document}
